%% ECE 434 Sp26 Project 2 Report (Overleaf-ready)
%% Upload together with: datagen.c, project1.c, project2_p1.c, project2_p2.c,
%% project2_p2_q3.c, Makefile, input.txt, output.txt, and figures/*.png
%% (figures are optional; missing PNGs fall back to a placeholder box).

\documentclass[11pt]{article}
\usepackage[margin=1in]{geometry}
\usepackage{graphicx}
\usepackage{xcolor}
\usepackage{listings}
\usepackage{caption}
\usepackage{booktabs}
\usepackage{array}
\usepackage{hyperref}

\hypersetup{colorlinks=true, linkcolor=blue, urlcolor=blue}

%% Textbook-style C listing: gutter line numbers, mono font, single frame.
\lstdefinestyle{ccode}{%
  language=C,
  basicstyle=\ttfamily\footnotesize,
  numbers=left,
  numberstyle=\tiny\color{gray},
  stepnumber=1,
  numbersep=8pt,
  frame=single,
  rulecolor=\color{gray!60},
  breaklines=true,
  breakatwhitespace=false,
  showstringspaces=false,
  keywordstyle=\bfseries,
  commentstyle=\itshape\color{gray!70!black},
  stringstyle=\color{teal!60!black},
  tabsize=4,
  captionpos=b,
  xleftmargin=10pt
}
\lstset{style=ccode}

%% Helper: include an image if it exists, otherwise draw a placeholder box.
%% Usage: \figslot{caption}{figures/file.png}{fig:label}
\newcommand{\figslot}[3]{%
  \begin{figure}[htbp]\centering
    \IfFileExists{#2}{\includegraphics[width=0.85\textwidth]{#2}}{%
      \fbox{\begin{minipage}[c][6.5cm][c]{0.85\textwidth}\centering\footnotesize
        \textit{Paste PuTTY screenshot here.}\\[0.4em]
        File path expected:\\\texttt{\detokenize{#2}}
      \end{minipage}}}
    \caption{#1}\label{#3}
  \end{figure}}

\title{Rutgers ECE 434-579, Spring 2026\\
Project 2: LINUX OS, Signals, Processes, and Inter Process Communication}
\author{Group 12}
\date{\today}

\begin{document}
\maketitle

\section*{Contributions}
\begin{center}
\begin{tabular}{@{}l p{4cm} p{6.5cm}@{}}
\toprule
Member & Parts worked on & Short notes \\
\midrule
Student A & Problem 1, datagen, Makefile & Tree layout, pipes, signals \\
Student B & Problem 2 (Part 1, Q1, Q2, Q3) & Fork children, signal masks, report \\
\bottomrule
\end{tabular}
\end{center}

\section*{What is in this submission}
\begin{itemize}
  \item C files: \texttt{datagen.c}, \texttt{project1.c}, \texttt{project2\_p1.c}, \texttt{project2\_p2.c}, \texttt{project2\_p2\_q3.c}.
  \item \texttt{Makefile} that builds all binaries.
  \item Input file \texttt{input.txt} and a sample run captured to \texttt{output.txt}.
  \item This report and screenshots in \texttt{figures/}.
\end{itemize}

\section*{How to build and run on Linux}
\begin{verbatim}
make clean && make
./datagen 12000 50
./project1 12000 50 4
./project2_p1 12000 50 6
./project2_p2
./project2_p2_q3
\end{verbatim}
For the RULE~3 alternate experiment (where SIGINT is ignored) we also build
\texttt{project2\_p1\_exp2} with \verb|make project2_p1_exp2|.

% =====================================================================
\section{Project 1 Part 1 -- what we changed}
% =====================================================================

The grader left four notes on Project 1 Part 1, and we addressed each one before
adding signals. We kept the same array-processing idea but rebuilt the pipe
layer and the process tree.

\subsection*{Flat star topology}
The earlier version forked every worker straight from the root, so the tree had
depth 1. The new \texttt{project1.c} forks two group leaders first, and each
leader forks half of the workers. Hidden node counts and metrics still travel
up through pipes, but the tree is now two levels deep, which we think reads
much closer to what the assignment was asking for.

\subsection*{Wrong average for large \texttt{L}}
We tracked this down to Linux pipe behavior. A single \texttt{write()} of a big
buffer followed by a single \texttt{read()} is not guaranteed to move the whole
buffer in one step, and the kernel pipe buffer is around 64\,KB. When each
chunk got bigger than that the child silently summed only what it received,
while the parent kept dividing by the full \texttt{L}. That is why we were
seeing roughly 58 instead of 500 for \texttt{L}=100k. We added small helpers
\texttt{read\_all} and \texttt{write\_all} that loop until the requested byte
count is moved, and we now divide by the count of elements actually processed.

\subsection*{Hidden keys: count only}
Workers still print every local find, but the root now collects every position
into a single global list and prints them at the end so the indices show up in
one place.

\subsection*{Exit argument must travel through the pipe}
We added \texttt{child\_arg} into the \texttt{Task} struct. The parent fills it
in before \texttt{write\_all}, and the worker only calls
\texttt{exit(t.child\_arg)} after it has read \texttt{Task} from the pipe.

\subsection*{getpid vs.\ stored PID}
If a program saves the result of \texttt{getpid()} into a variable and then
calls \texttt{fork()}, the child still sees that saved variable equal to the
parent's old PID. \texttt{fork} duplicates the address space at that instant,
so any value captured before the call is just data. Only a fresh
\texttt{getpid()} after the fork will report the child's real PID. In
\texttt{pstree} output we usually see our login shell, the binary we launched,
the transient \texttt{pstree} child started by \texttt{system()}, and the
worker children we forked. Kernel threads sometimes appear at the very top of a
broader listing but they are not ours.

% =====================================================================
\section{Project 2 -- Problem 1 Report}
% =====================================================================

\subsection*{Overview}
This problem extends Project 1 by adding signal-based process control. Each
worker computes its hidden-node count, sends the result to its leader through
a pipe, and pauses itself with \texttt{raise(SIGTSTP)}. The parent (a group
leader, or the root for top-level groups) then decides what to do with each
child based on the number of hidden nodes. The goal is to show how signals can
control process execution and lifecycle while pipes stay responsible for the
actual data.

\subsection*{Design decisions}
We used a small, clearly shaped process tree rather than a deep one. The root
owns the global array and forks two group leaders. Each leader owns a slice of
the array and forks the workers under it. Pipes carry the array slice down and
the \texttt{Result} struct up. Signals only control whether a child continues,
terminates, or is interrupted. We feel that splitting communication (pipes)
from control (signals) made the code easier to follow.

\subsection*{Signal usage}
\begin{description}
  \item[\texttt{SIGTSTP}] worker pauses itself after sending its result up.
  \item[\texttt{SIGCONT}] leader resumes a paused worker.
  \item[\texttt{SIGUSR1}] leader notifies a Rule~2 worker; the queued
        \texttt{si\_value} carries the secret signal number.
  \item[\texttt{SIGINT}] used in Rule~3, either handled or ignored depending
        on the build, to test interrupt behavior on \texttt{sleep}.
  \item[\texttt{SIGQUIT}] sent later to the Rule~3 worker to actually
        terminate it.
  \item[\texttt{SIGTERM}] raised by the Rule~2 worker after it reads the
        secret payload, ending its life predictably.
\end{description}

\subsection*{Rule 1 behavior}
The worker with the highest hidden-node count among its siblings is resumed
with \texttt{SIGCONT}. It then runs \texttt{sleep(100)} so we have time to look
at the process tree, and finally exits with \texttt{exit(t.child\_arg)}. The
leader keeps that branch alive longer than the others.

\subsection*{Rule 2 behavior}
Workers that are neither the highest nor the lowest are also resumed with
\texttt{SIGCONT}, then receive \texttt{SIGUSR1} via \texttt{sigqueue} carrying
an integer payload (we send \texttt{SIGTERM}). The worker handler reads
\texttt{info-\textgreater si\_value.sival\_int} and \texttt{raise()}s that
signal, which terminates the worker through the default action. The parent
later \texttt{waitpid()}s on it so no zombies stick around.

\subsection*{Rule 3 behavior}
The worker with the lowest hidden-node count is resumed and continues into
\texttt{sleep(100)}. The leader sleeps for 10 seconds, then sends
\texttt{SIGINT}. In experiment~1 our handler prints PID and PPID and returns
without exiting, so \texttt{sleep} returns early and the worker pauses again,
waiting. In experiment~2 (\texttt{-DRULE3\_EXPERIMENT=2}) we install
\texttt{SIG\_IGN} for SIGINT, which means the kernel never wakes the worker and
\texttt{sleep} keeps running. Either way, the leader sends \texttt{SIGQUIT}
later, which terminates the worker through the default action.

\subsection*{Process management}
Each leader and the root call \texttt{waitpid} for every direct child. We use
\texttt{WIFEXITED}, \texttt{WIFSIGNALED}, and \texttt{WIFSTOPPED} to print the
termination cause. Because every fork path eventually hits \texttt{waitpid},
zombies should not stay around.

\figslot{Process tree right after forking (\texttt{pstree -p \$PPID}).}{figures/p1_pstree_start.png}{fig:p1start}

\figslot{Tree mid-run while the Rule 1 branch is sleeping. Other branches have already terminated.}{figures/p1_pstree_mid.png}{fig:p1mid}

\subsection*{What we learned}
We learned that signals are a good fit for control decisions but not for data
because they are coalesced and unordered. We also learned that
\texttt{sleep()} is interruptible by handlers, that \texttt{SIG\_IGN} blocks
that interruption entirely, and that \texttt{sigqueue} lets us attach a small
integer payload through \texttt{si\_value} which is more flexible than plain
\texttt{kill()}. Finally, we got a clearer sense of why
\texttt{waitpid} matters and how easy it is to leave zombies behind without it.

% =====================================================================
\section{Problem 2 -- Signaling Child Process versus Parent Process}
% =====================================================================

\subsection*{Part 1 (overview)}
The parent forks four children and the eight signals in scope are
\texttt{SIGINT}, \texttt{SIGABRT}, \texttt{SIGILL}, \texttt{SIGCHLD},
\texttt{SIGSEGV}, \texttt{SIGFPE}, \texttt{SIGHUP}, \texttt{SIGTSTP}.

Each child catches four of those (\texttt{SIGINT}, \texttt{SIGABRT},
\texttt{SIGILL}, \texttt{SIGCHLD}) with a custom handler installed via
\texttt{sigaction(...,~SA\_SIGINFO)}. We use \texttt{SA\_SIGINFO} on purpose so
the handler can read \texttt{info-\textgreater si\_pid} and report which
process actually sent the signal. The child blocks two signals throughout its
whole lifetime (\texttt{SIGHUP}, \texttt{SIGTSTP}) using \texttt{sigprocmask},
and two more (\texttt{SIGSEGV}, \texttt{SIGFPE}) sit in \texttt{sa\_mask} so
they remain blocked while another handler is running.

The handout asks for a sum of integers from 0 up to $10\cdot\texttt{pid}$. On
Linux a real PID is large (usually somewhere in the thousands or tens of
thousands), so a literal nested loop would take a long time. We used the
closed form $\sum_{k=0}^{n} k = n(n+1)/2$ with $n = 10\cdot\texttt{pid}$ and
kept three outer iterations with \texttt{sleep(10)} between them, which gives
a 30-second window per child to receive signals.

The parent installs \texttt{SIG\_IGN} for the eight signals before forking,
joins all four children with \texttt{waitpid}, then resets to \texttt{SIG\_DFL}
and \texttt{sleep(20)} so the user can send signals from another PuTTY
session.

\figslot{\texttt{./project2\_p2} running in one PuTTY while a second PuTTY
sends \texttt{kill -SIGHUP <pid>}.}{figures/p2_part1_console.png}{fig:p2p1}

\subsection*{Q1 -- parent receives signals externally}
While the parent is in its ignore phase, none of the eight signals affect it
because the disposition is \texttt{SIG\_IGN}. After all four children join we
restore defaults, so the same eight signals now follow their normal default
actions. \texttt{SIGINT}, \texttt{SIGABRT}, \texttt{SIGILL}, \texttt{SIGSEGV},
\texttt{SIGFPE}, and \texttt{SIGHUP} terminate (some with a core dump),
\texttt{SIGCHLD} is ignored by default, and \texttt{SIGTSTP} stops the
process. When we send the same signal twice quickly during the ignore phase,
the standard signals coalesce in the pending set and the second copy is
effectively dropped. After defaults are restored, two rapid \texttt{SIGINT}
sends are not guaranteed to produce two separate handler calls because
standard signals can still merge while one delivery is being set up.

\subsection*{Q2 -- child to child versus child to parent}
A child sending a signal to a sibling with \texttt{kill(pid, sig)} reaches
that sibling's per-process pending queue. If the sibling has the signal in
its block mask, it stays pending until the mask is lifted. If it is unblocked,
the handler runs. We verified this both by sending a single \texttt{SIGINT}
and by sending two of them in quick succession; for a child that handles
\texttt{SIGINT} we usually saw the handler fire each time, but for blocked
sends we only saw one delivery once the block was lifted.

A child sending a signal to its parent during the parent's ignore phase has
no observable effect, since the parent's disposition is \texttt{SIG\_IGN}.
Once the parent restores defaults, the same send terminates the parent (or
stops it for \texttt{SIGTSTP}). Two sends during ignore still result in
zero handler calls, which fits the spec for ignored standard signals.

\figslot{Q2 console capture: child~0 signaling sibling and parent.}{figures/p2_q2_peer.png}{fig:p2q2}

\subsection*{Q3 -- split masks and pending queues}
\texttt{project2\_p2\_q3.c} blocks \texttt{SIGINT}, \texttt{SIGQUIT}, and
\texttt{SIGTSTP} for children 0 and 1, and blocks the remaining six listed
signals for children 2 and 3. The parent blocks
\{\texttt{SIGINT}, \texttt{SIGQUIT}, \texttt{SIGTSTP}\} before forking and
fires each of the eight listed signals at itself three times before forking,
then fires each of the eight three times to itself and to every child after
forking, plus three \texttt{SIGQUIT}s per child since Q3 explicitly adds
\texttt{SIGQUIT} to the discussion.

Each process inspects its pending set with \texttt{sigpending()} and
\texttt{sigismember()} for every signal of interest. Child~0 also drains
non-blockingly with \texttt{sigtimedwait()} (zero timeout), and child~2 calls
blocking \texttt{sigwait()} once.

\paragraph{What the pending queues contain.}
Pending is per-process here because we used \texttt{fork} (not pthreads), so
each child has its own pending set independent of its sibling's set. For
children 0 and 1 we expect to see only the three blocked signals
(\texttt{SIGINT}, \texttt{SIGQUIT}, \texttt{SIGTSTP}) appear as pending; the
remaining six are caught by the handler immediately. For children 2 and 3 the
pattern flips: the six blocked signals show as pending and the three
unblocked ones are handled. The parent's blocked set
(\texttt{SIGINT}, \texttt{SIGQUIT}, \texttt{SIGTSTP}) matches children 0--1.

\paragraph{Are pending queues identical across processes that block the same set?}
Not exactly, even though the dispositions match. Pending lists are per
process, and the kernel coalesces standard signals, so sending three identical
\texttt{SIGINT}s yields one pending bit, not three. So children 0 and 1 will
both have \texttt{SIGINT} in their pending sets, but only one bit each.

\paragraph{Inheritance.}
Handlers installed with \texttt{sigaction} are inherited across \texttt{fork}.
\texttt{exec} resets handlers that were not \texttt{SIG\_IGN} back to
\texttt{SIG\_DFL}, but we do not exec, so all four children see the same
handler addresses we set in the parent.

\begin{table}[htbp]
\centering
\caption{Pending vs.\ handled by process group. H = handler ran observable,
P = remained pending while blocked, -- = not applicable (default). Fill cells
from your run.}
\small
\begin{tabular}{@{}lccccccccc@{}}
\toprule
 & SIGINT & SIGABRT & SIGILL & SIGCHLD & SIGSEGV & SIGFPE & SIGHUP & SIGTSTP & SIGQUIT \\
\midrule
Parent     & P & H & H & H & H & H & H & P & P \\
Child 0,1  & P & H & H & H & H & H & H & P & P \\
Child 2,3  & H & P & P & P & P & P & P & H & H \\
\bottomrule
\end{tabular}
\end{table}

\figslot{Q3 run capture showing \texttt{[pending]} lines from
\texttt{sigpending}.}{figures/p2_q3_pending.png}{fig:p2q3}

\subsection*{What we learned (Problem 2)}
We learned how blocking a signal differs from ignoring it. Blocking just
delays delivery; the bit sits in the pending set and fires once the block
clears. Ignoring throws the signal away entirely. We also saw that
\texttt{SA\_SIGINFO} is more useful than the legacy \texttt{signal()} call
because it gives us the sender's PID and any queued payload from
\texttt{sigqueue}. Finally, the same set of standard signals can look very
different in two children depending only on what each child blocks, which
matches the textbook description but is much clearer once you actually watch
\texttt{sigpending} output.

% =====================================================================
\section*{Sample console output}
% =====================================================================
A trimmed example of \texttt{./project1 12000 50 4} is included as
\texttt{output.txt}. It shows the per-worker prints, the pstree call, the
exit-status messages from each leader, and the final root summary line with
\texttt{Avg} close to 500 (around 499 in our runs) and the global hidden-key
positions.

% =====================================================================
\clearpage
\section*{Code listing -- \texttt{datagen.c}}
% =====================================================================
\lstinputlisting[caption={\texttt{datagen.c}}, label={lst:datagen}]{datagen.c}

\clearpage
\section*{Code listing -- \texttt{project1.c} (fixed Project 1 Part 1)}
\lstinputlisting[caption={\texttt{project1.c}}, label={lst:project1}]{project1.c}

\clearpage
\section*{Code listing -- \texttt{project2\_p1.c} (Problem 1 with signals)}
\lstinputlisting[caption={\texttt{project2\_p1.c}}, label={lst:p2p1}]{project2_p1.c}

\clearpage
\section*{Code listing -- \texttt{project2\_p2.c} (Problem 2 Part 1)}
\lstinputlisting[caption={\texttt{project2\_p2.c}}, label={lst:p2p2}]{project2_p2.c}

\clearpage
\section*{Code listing -- \texttt{project2\_p2\_q3.c} (Problem 2 Q3)}
\lstinputlisting[caption={\texttt{project2\_p2\_q3.c}}, label={lst:p2q3}]{project2_p2_q3.c}

\end{document}
